% hybrid_ar_ctc_and_tying_ablation.tex
% Thesis-ready text: Hybrid AR+CTC training and the weight-tying ablation.
% Drop into your thesis with: % hybrid_ar_ctc_and_tying_ablation.tex
% Thesis-ready text: Hybrid AR+CTC training and the weight-tying ablation.
% Drop into your thesis with: % hybrid_ar_ctc_and_tying_ablation.tex
% Thesis-ready text: Hybrid AR+CTC training and the weight-tying ablation.
% Drop into your thesis with: % hybrid_ar_ctc_and_tying_ablation.tex
% Thesis-ready text: Hybrid AR+CTC training and the weight-tying ablation.
% Drop into your thesis with: \input{<path>/hybrid_ar_ctc_and_tying_ablation.tex}

\section{Hybrid autoregressive--CTC training}
\label{sec:hybrid-ar-ctc}

\paragraph{Motivation.}
Autoregressive (AR) sequence-to-sequence decoders with cross-attention are strong recognizers, but training can benefit from an auxiliary objective that encourages monotonic, frame-synchronous alignment.
Connectionist Temporal Classification (CTC) provides such a signal by defining a differentiable alignment marginalization over encoder time steps.
In this work, we adopt a \emph{hybrid} training regime where an AR decoder is the primary recognizer, while an auxiliary CTC head regularizes the shared encoder representation.

\subsection{Model}
Let $\mathbf{x} \in \mathbb{R}^{T \times C}$ denote an input IMU sequence of length $T$ with $C$ sensor channels.
An encoder $f_{\text{enc}}$ produces a downsampled sequence of hidden states
\begin{equation}
\mathbf{h} = f_{\text{enc}}(\mathbf{x}) \in \mathbb{R}^{T' \times d_{\text{enc}}},
\end{equation}
where $T' \le T$ due to temporal downsampling.
An AR decoder $f_{\text{ar}}$ attends to $\mathbf{h}$ and predicts a token sequence $\mathbf{y} = (y_1,\dots,y_N)$ via teacher forcing during training.

\paragraph{Vocabularies and special tokens.}
We use a character-level vocabulary (``categories'') with the CTC blank symbol at index $0$.
For AR decoding we use the same base vocabulary and append three special symbols \texttt{PAD}, \texttt{BOS}, \texttt{EOS} at the end.
Thus, for character mode,
\begin{equation}
V_{\text{ctc}} = |\mathcal{V}_{\text{ctc}}|, \qquad
V_{\text{ar}} = V_{\text{ctc}} + 3.
\end{equation}
This layout ensures that token IDs $0,\dots,V_{\text{ctc}}-1$ have the same semantics for both CTC and AR.

\subsection{Objectives}
\paragraph{AR cross-entropy.}
With teacher forcing, the decoder defines next-token probabilities
\begin{equation}
P_{\text{ar}}(y_t \mid y_{<t}, \mathbf{h}), \qquad t=1,\dots,N.
\end{equation}
The AR loss is the standard cross-entropy (optionally with label smoothing) summed over positions:
\begin{equation}
\mathcal{L}_{\text{ar}} = - \sum_{t=1}^{N} \log P_{\text{ar}}(y_t \mid y_{<t}, \mathbf{h}).
\end{equation}

\paragraph{CTC loss.}
A per-timestep classifier produces logits over $\mathcal{V}_{\text{ctc}}$ for each encoder frame.
CTC defines the conditional likelihood by summing over all monotonic alignments that collapse (via blank removal and repetition merging) to the target sequence.
We use the negative log-likelihood form
\begin{equation}
\mathcal{L}_{\text{ctc}} = -\log P_{\text{ctc}}(\mathbf{y} \mid \mathbf{h}).
\end{equation}

\paragraph{Hybrid loss.}
The overall hybrid loss combines both terms with weights $\lambda_{\text{ar}}$ and $\lambda_{\text{ctc}}$:
\begin{equation}
\mathcal{L}_{\text{hyb}} = \lambda_{\text{ar}}\, \mathcal{L}_{\text{ar}} + \lambda_{\text{ctc}}\, \mathcal{L}_{\text{ctc}}.
\label{eq:hybrid-loss}
\end{equation}
In our experiments, $\lambda_{\text{ar}}=1$ and $\lambda_{\text{ctc}}$ is treated as a hyperparameter (optionally scheduled across epochs).

\paragraph{Why this helps.}
The CTC term provides an auxiliary alignment/regularization signal on the encoder time axis.
This can stabilize optimization and improve generalization, especially in settings with noisy sensor sequences and variable writing styles.
Importantly, the AR decoder remains the primary recognizer at inference time.

\section{Ablation: ``single-decoder-ish'' CTC via weight tying}
\label{sec:ctc-weight-tying}

\subsection{Goal}
A common critique of hybrid models is that gains might come from adding extra parameters (the auxiliary head) rather than from the auxiliary objective.
To control for this, we introduce an ablation where the CTC classifier reuses the AR decoder's output projection weights.
This produces a more ``single-decoder-ish'' setup: CTC still scores encoder frames, but it does not introduce a separate token classifier matrix.

\subsection{Baseline (separate CTC head)}
In the standard hybrid model, CTC logits are computed via a dedicated linear layer
\begin{equation}
\mathbf{z}^{\text{ctc}}_t = \mathbf{W}_{\text{ctc}}\, \mathbf{h}_t + \mathbf{b}_{\text{ctc}},
\qquad \mathbf{W}_{\text{ctc}} \in \mathbb{R}^{V_{\text{ctc}} \times d_{\text{enc}}}.
\end{equation}
This adds $V_{\text{ctc}}\cdot d_{\text{enc}}$ parameters (plus bias).

\subsection{Variant A (primary ablation): tie CTC to the AR output projection}
\paragraph{AR output projection.}
The AR decoder produces hidden states in decoder space $\mathbf{s}_t \in \mathbb{R}^{d_{\text{dec}}}$ and maps them to vocabulary logits using its output projection
\begin{equation}
\mathbf{z}^{\text{ar}}_t = \mathbf{W}_{\text{ar}}\, \mathbf{s}_t + \mathbf{b}_{\text{ar}},
\qquad \mathbf{W}_{\text{ar}} \in \mathbb{R}^{V_{\text{ar}} \times d_{\text{dec}}}.
\end{equation}

\paragraph{CTC logits from encoder frames using shared classifier rows.}
To reuse the same classifier weights, we compute CTC logits by taking the \emph{prefix slice} of the AR projection matrix corresponding to the shared vocabulary IDs:
\begin{equation}
\mathbf{W}_{\text{ctc}} \equiv \mathbf{W}_{\text{ar}}[0{:}V_{\text{ctc}}],
\qquad \mathbf{b}_{\text{ctc}} \equiv \mathbf{b}_{\text{ar}}[0{:}V_{\text{ctc}}].
\end{equation}
Since CTC operates on encoder frames, we must ensure its input feature dimension matches $d_{\text{dec}}$.
When $d_{\text{enc}} \ne d_{\text{dec}}$, we reuse the existing encoder-to-decoder projection (already required for AR cross-attention in our implementation):
\begin{equation}
\tilde{\mathbf{h}}_t = \mathbf{P}\, \mathbf{h}_t, \qquad \mathbf{P} \in \mathbb{R}^{d_{\text{dec}} \times d_{\text{enc}}}.
\end{equation}
The tied CTC logits become
\begin{equation}
\mathbf{z}^{\text{ctc}}_t = \mathbf{W}_{\text{ar}}[0{:}V_{\text{ctc}}] \, \tilde{\mathbf{h}}_t + \mathbf{b}_{\text{ar}}[0{:}V_{\text{ctc}}].
\label{eq:tied-ctc-logits}
\end{equation}
With this design, there is no independent $\mathbf{W}_{\text{ctc}}$ parameter matrix.
Gradients from the CTC loss update the same token classifier used by the AR decoder (for the shared token subset).

\paragraph{Effect on parameter count and reporting.}
Compared to the baseline hybrid model with a separate linear CTC head, this ablation removes the dedicated CTC classifier parameters $\mathbf{W}_{\text{ctc}} \in \mathbb{R}^{V_{\text{ctc}}\times d_{\text{enc}}}$ (and optional bias), i.e., it eliminates approximately $V_{\text{ctc}}\cdot d_{\text{enc}}$ learnable weights.
No new classifier parameters are introduced because CTC reuses the AR output projection rows.
When $d_{\text{enc}}\ne d_{\text{dec}}$, we reuse the existing encoder-to-decoder projection already required for cross-attention (thus the ablation does not add an extra ``adapter'' beyond the baseline AR model).
In the thesis, we report parameter counts and MACs for each variant in a single table (baseline hybrid vs tied-CTC), and explicitly attribute any reduction in parameters to the removal of the standalone CTC head.

\paragraph{Interpretation.}
This ablation isolates whether improvements from hybrid training are attributable to the auxiliary alignment objective (CTC) rather than increased head capacity.
It also aligns with the intuition that the decoder is the primary recognizer, while CTC provides complementary supervision using the same classifier.

\subsection{Configuration switches for the ablation}
We expose minimal switches in the configuration under \texttt{dual\_head.tie}:
\begin{itemize}
  \item \texttt{ctc\_to\_ar\_outproj}: enables the tied CTC classifier (Variant A).
  \item \texttt{ar\_emb\_outproj}: optional secondary ablation that ties the AR embedding matrix to the AR output projection (LM-style weight tying).
  \item \texttt{ctc\_input\_space}: \texttt{"dec"} (recommended) uses the existing encoder$\to$decoder projection so that Eq.~\eqref{eq:tied-ctc-logits} is well-defined even when $d_{\text{enc}}\ne d_{\text{dec}}$.
\end{itemize}
A concrete YAML snippet is:
\begin{verbatim}
dual_head:
  enabled: true
  arch_ctc: linear
  lambda_ar: 1.0
  lambda_ctc: 0.6
  tie:
    ctc_to_ar_outproj: true
    ar_emb_outproj: false
    ctc_input_space: "dec"
\end{verbatim}

\subsection{Optional auxiliary supervision inside the decoder}
\label{sec:decoder-aux-supervision}
To more directly regularize the AR decoder (``between decoder layers''), we additionally implement two optional training-time losses.
These ablations are available both in hybrid AR+CTC mode (under \texttt{dual\_head.*}) and in AR-only mode (under top-level config keys), enabling fair comparisons.

\paragraph{Deep supervision (auxiliary token cross-entropy at intermediate layers).}
When enabled, we compute token-level cross-entropy not only from the final decoder layer but also from intermediate layers (e.g., layers 1--3 for a 4-layer decoder).
Crucially, all layers reuse the \emph{same} vocabulary projection matrix (the AR output projection), so this adds no independent classifier capacity.

For hybrid AR+CTC mode, the configuration lives under \texttt{dual\_head.deep\_supervision}:
\begin{verbatim}
dual_head:
  deep_supervision:
    enabled: true
    layers: [1, 2, 3]      # supervised intermediate layers (1-indexed)
    include_final: true    # also supervise the final layer
    reduce: mean           # mean|sum across supervised layers
\end{verbatim}

For AR-only mode, use a top-level block:
\begin{verbatim}
deep_supervision:
  enabled: true
  layers: [1, 2, 3]
  include_final: true
  reduce: mean
\end{verbatim}

\paragraph{Decoder-side ``CTC-like'' loss (CTC on encoder frames conditioned on decoder states).}
True CTC is defined over a frame-synchronous time axis, whereas decoder layers are naturally defined over the token axis.
To obtain a decoder-internal auxiliary signal while remaining mathematically consistent with CTC, we construct frame-length ($T'$) logits conditioned on intermediate decoder states.
The token classifier itself is still shared (we reuse the AR output projection rows for the shared vocabulary prefix), so we do not introduce an additional classifier matrix.

This option implements the suggestion to ``put the CTC head with the transformer decoder as a regularizer'' (see discussion with supervisor).

Two modes are available via \texttt{decoder\_side\_ctc.mode}:
\begin{itemize}
  \item \textbf{attention} (default): Encoder frames (queries) attend to decoder token states (keys/values) via a small reverse cross-attention module. This adds $\sim 3 d^2$ parameters (Q/K/V projections + output) but allows each frame to selectively attend to relevant decoder positions.
  \item \textbf{pool}: Global average pooling of decoder states, then broadcast-added to encoder frames. This adds \emph{zero} extra parameters and provides a simpler baseline for testing decoder-conditioning.
\end{itemize}

For hybrid AR+CTC mode, the configuration is under \texttt{dual\_head.decoder\_side\_ctc}:
\begin{verbatim}
dual_head:
  decoder_side_ctc:
    enabled: true
    mode: attention      # "attention" or "pool"
    layers: [1, 2, 3]    # which decoder layer states condition the frame logits
    lambda: 0.1          # weight of this auxiliary loss
    nhead: 4             # attention heads (only used if mode=attention)
    layernorm: true      # layer norm after attention (only if mode=attention)
\end{verbatim}

For AR-only mode (AR + decoder-side CTC ablation), use a top-level block:
\begin{verbatim}
decoder_side_ctc:
  enabled: true
  mode: pool             # zero extra params
  layers: [1, 2, 3]
  lambda: 0.1
\end{verbatim}

\paragraph{Ablation matrix.}
With these switches, the following conditions can be compared on a common encoder/decoder backbone:
\begin{enumerate}
  \item \textbf{AR-only}: pure autoregressive CE (baseline).
  \item \textbf{AR + deep supervision}: CE at layers 1/2/3/4 (no extra head capacity).
  \item \textbf{AR + decoder-side CTC}: CTC regularization ``between decoder layers'' (adds reverse-attention params only).
  \item \textbf{Hybrid AR+CTC (encoder-side)}: separate CTC head after encoder (baseline hybrid).
  \item \textbf{Hybrid AR+CTC + deep supervision / decoder-side CTC}: combinations of the above.
\end{enumerate}
This design directly tests whether placing auxiliary supervision closer to the decoder (the professor's suggestion) improves recognition compared to placing it after the encoder.

\paragraph{Secondary ablation and experimental design.}
The switch \texttt{ar\_emb\_outproj=true} enables AR embedding--classifier tying (i.e., the same weight matrix is used for token lookup and for the final vocabulary projection).
This can change optimization dynamics even without CTC weight sharing.
For clean attribution, we therefore treat it as an \emph{independent} ablation and report it both (i) on its own and (ii) in combination with Variant~A.
Concretely, we run a small factorial comparison: baseline hybrid (separate CTC head), Variant~A (CTC$\leftrightarrow$out\_proj tied), AR-only tying (\texttt{ar\_emb\_outproj=true} only), and ``full tying'' (both flags enabled).

\subsection{Results and conclusion (capacity-control test)}
\label{sec:ctc-weight-tying-results}
Table~\ref{tab:hybrid-ctc-tying-results} reports the word-level results of the AR-only baseline and hybrid AR+CTC training under different tying configurations.
All numbers are character error rate (CER) and word error rate (WER) in \%, using the same encoder/decoder and evaluation protocol across runs.
We sweep the CTC weight $\lambda_{\text{ctc}}\in\{0.5,0.6\}$, keeping $\lambda_{\text{ar}}=1$.

\begin{table}[t]
\centering
\begin{tabular}{lcc}
\hline
\textbf{Variant} & \textbf{CER (\%)} & \textbf{WER (\%)} \\
\hline
AR-only (no CTC) & 6.82 & 10.31 \\
\hline
Hybrid AR+CTC (separate CTC head), $\lambda_{\text{ctc}}=0.6$ & 6.46 & 9.76 \\
Hybrid AR+CTC (separate CTC head), $\lambda_{\text{ctc}}=0.5$ & 6.48 & 9.87 \\
\hline
Hybrid + tie \texttt{ar\_emb\_outproj} only, $\lambda_{\text{ctc}}=0.6$ & 6.53 & 9.97 \\
Hybrid + tie \texttt{ar\_emb\_outproj} only, $\lambda_{\text{ctc}}=0.5$ & 6.69 & 10.12 \\
\hline
Hybrid + tie \texttt{ctc\_to\_ar\_outproj} only, $\lambda_{\text{ctc}}=0.6$ & 7.00 & 10.49 \\
Hybrid + tie \texttt{ctc\_to\_ar\_outproj} only, $\lambda_{\text{ctc}}=0.5$ & 6.99 & 10.53 \\
\hline
Hybrid + full tying (both flags), $\lambda_{\text{ctc}}=0.5$ & 7.22 & 11.32 \\
\hline
\end{tabular}
\caption{Ablation of hybrid AR+CTC training and weight tying (word-level recognition). ``Separate CTC head'' denotes the standard hybrid model with an independent CTC classifier matrix $\mathbf{W}_{\text{ctc}}$. ``Tie \texttt{ctc\_to\_ar\_outproj}'' removes this independent classifier by reusing the AR output projection rows as in Eq.~\eqref{eq:tied-ctc-logits}.}
\label{tab:hybrid-ctc-tying-results}
\end{table}

\paragraph{What the capacity-control ablation tells us.}
The untied hybrid model (separate CTC head) improves over AR-only for both $\lambda_{\text{ctc}}=0.5$ and $0.6$.
However, when we apply the intended capacity-control by tying the CTC classifier to the AR output projection (\texttt{ctc\_to\_ar\_outproj=true}), performance drops back to (and slightly below) the AR-only baseline.
Therefore, based on these experiments, we \emph{cannot} conclude that the hybrid gains come purely from the auxiliary objective independent of added head capacity.
The results are consistent with two (not mutually exclusive) explanations: (i) the standalone CTC head provides useful classifier flexibility beyond the shared AR projection, and/or (ii) the tying constraint introduces optimization trade-offs (gradient interference) that negate the hybrid benefit.

\paragraph{Secondary observation (embedding/output tying).}
Tying the AR embedding to the AR output projection (\texttt{ar\_emb\_outproj=true}) does not improve performance in this setting and slightly degrades CER/WER compared to the untied hybrid.
This reinforces the importance of treating embedding/output tying as an independent ablation rather than as a ``free'' control.

\subsection{Reproducibility and sanity checks}
When tying is enabled, we enforce the following invariants at model initialization:
\begin{itemize}
  \item $V_{\text{ar}} \ge V_{\text{ctc}}$ so that a prefix slice exists.
  \item Character-mode special token layout: if $V_{\text{ar}} = V_{\text{ctc}} + 3$, then \texttt{PAD}$=V_{\text{ctc}}$, \texttt{BOS}$=V_{\text{ctc}}+1$, \texttt{EOS}$=V_{\text{ctc}}+2$.
  \item \texttt{ctc\_input\_space="enc"} is only allowed when $d_{\text{enc}}=d_{\text{dec}}$.
\end{itemize}
These checks ensure the tied classifier uses consistent token semantics across CTC and AR.

% End of file


\section{Hybrid autoregressive--CTC training}
\label{sec:hybrid-ar-ctc}

\paragraph{Motivation.}
Autoregressive (AR) sequence-to-sequence decoders with cross-attention are strong recognizers, but training can benefit from an auxiliary objective that encourages monotonic, frame-synchronous alignment.
Connectionist Temporal Classification (CTC) provides such a signal by defining a differentiable alignment marginalization over encoder time steps.
In this work, we adopt a \emph{hybrid} training regime where an AR decoder is the primary recognizer, while an auxiliary CTC head regularizes the shared encoder representation.

\subsection{Model}
Let $\mathbf{x} \in \mathbb{R}^{T \times C}$ denote an input IMU sequence of length $T$ with $C$ sensor channels.
An encoder $f_{\text{enc}}$ produces a downsampled sequence of hidden states
\begin{equation}
\mathbf{h} = f_{\text{enc}}(\mathbf{x}) \in \mathbb{R}^{T' \times d_{\text{enc}}},
\end{equation}
where $T' \le T$ due to temporal downsampling.
An AR decoder $f_{\text{ar}}$ attends to $\mathbf{h}$ and predicts a token sequence $\mathbf{y} = (y_1,\dots,y_N)$ via teacher forcing during training.

\paragraph{Vocabularies and special tokens.}
We use a character-level vocabulary (``categories'') with the CTC blank symbol at index $0$.
For AR decoding we use the same base vocabulary and append three special symbols \texttt{PAD}, \texttt{BOS}, \texttt{EOS} at the end.
Thus, for character mode,
\begin{equation}
V_{\text{ctc}} = |\mathcal{V}_{\text{ctc}}|, \qquad
V_{\text{ar}} = V_{\text{ctc}} + 3.
\end{equation}
This layout ensures that token IDs $0,\dots,V_{\text{ctc}}-1$ have the same semantics for both CTC and AR.

\subsection{Objectives}
\paragraph{AR cross-entropy.}
With teacher forcing, the decoder defines next-token probabilities
\begin{equation}
P_{\text{ar}}(y_t \mid y_{<t}, \mathbf{h}), \qquad t=1,\dots,N.
\end{equation}
The AR loss is the standard cross-entropy (optionally with label smoothing) summed over positions:
\begin{equation}
\mathcal{L}_{\text{ar}} = - \sum_{t=1}^{N} \log P_{\text{ar}}(y_t \mid y_{<t}, \mathbf{h}).
\end{equation}

\paragraph{CTC loss.}
A per-timestep classifier produces logits over $\mathcal{V}_{\text{ctc}}$ for each encoder frame.
CTC defines the conditional likelihood by summing over all monotonic alignments that collapse (via blank removal and repetition merging) to the target sequence.
We use the negative log-likelihood form
\begin{equation}
\mathcal{L}_{\text{ctc}} = -\log P_{\text{ctc}}(\mathbf{y} \mid \mathbf{h}).
\end{equation}

\paragraph{Hybrid loss.}
The overall hybrid loss combines both terms with weights $\lambda_{\text{ar}}$ and $\lambda_{\text{ctc}}$:
\begin{equation}
\mathcal{L}_{\text{hyb}} = \lambda_{\text{ar}}\, \mathcal{L}_{\text{ar}} + \lambda_{\text{ctc}}\, \mathcal{L}_{\text{ctc}}.
\label{eq:hybrid-loss}
\end{equation}
In our experiments, $\lambda_{\text{ar}}=1$ and $\lambda_{\text{ctc}}$ is treated as a hyperparameter (optionally scheduled across epochs).

\paragraph{Why this helps.}
The CTC term provides an auxiliary alignment/regularization signal on the encoder time axis.
This can stabilize optimization and improve generalization, especially in settings with noisy sensor sequences and variable writing styles.
Importantly, the AR decoder remains the primary recognizer at inference time.

\section{Ablation: ``single-decoder-ish'' CTC via weight tying}
\label{sec:ctc-weight-tying}

\subsection{Goal}
A common critique of hybrid models is that gains might come from adding extra parameters (the auxiliary head) rather than from the auxiliary objective.
To control for this, we introduce an ablation where the CTC classifier reuses the AR decoder's output projection weights.
This produces a more ``single-decoder-ish'' setup: CTC still scores encoder frames, but it does not introduce a separate token classifier matrix.

\subsection{Baseline (separate CTC head)}
In the standard hybrid model, CTC logits are computed via a dedicated linear layer
\begin{equation}
\mathbf{z}^{\text{ctc}}_t = \mathbf{W}_{\text{ctc}}\, \mathbf{h}_t + \mathbf{b}_{\text{ctc}},
\qquad \mathbf{W}_{\text{ctc}} \in \mathbb{R}^{V_{\text{ctc}} \times d_{\text{enc}}}.
\end{equation}
This adds $V_{\text{ctc}}\cdot d_{\text{enc}}$ parameters (plus bias).

\subsection{Variant A (primary ablation): tie CTC to the AR output projection}
\paragraph{AR output projection.}
The AR decoder produces hidden states in decoder space $\mathbf{s}_t \in \mathbb{R}^{d_{\text{dec}}}$ and maps them to vocabulary logits using its output projection
\begin{equation}
\mathbf{z}^{\text{ar}}_t = \mathbf{W}_{\text{ar}}\, \mathbf{s}_t + \mathbf{b}_{\text{ar}},
\qquad \mathbf{W}_{\text{ar}} \in \mathbb{R}^{V_{\text{ar}} \times d_{\text{dec}}}.
\end{equation}

\paragraph{CTC logits from encoder frames using shared classifier rows.}
To reuse the same classifier weights, we compute CTC logits by taking the \emph{prefix slice} of the AR projection matrix corresponding to the shared vocabulary IDs:
\begin{equation}
\mathbf{W}_{\text{ctc}} \equiv \mathbf{W}_{\text{ar}}[0{:}V_{\text{ctc}}],
\qquad \mathbf{b}_{\text{ctc}} \equiv \mathbf{b}_{\text{ar}}[0{:}V_{\text{ctc}}].
\end{equation}
Since CTC operates on encoder frames, we must ensure its input feature dimension matches $d_{\text{dec}}$.
When $d_{\text{enc}} \ne d_{\text{dec}}$, we reuse the existing encoder-to-decoder projection (already required for AR cross-attention in our implementation):
\begin{equation}
\tilde{\mathbf{h}}_t = \mathbf{P}\, \mathbf{h}_t, \qquad \mathbf{P} \in \mathbb{R}^{d_{\text{dec}} \times d_{\text{enc}}}.
\end{equation}
The tied CTC logits become
\begin{equation}
\mathbf{z}^{\text{ctc}}_t = \mathbf{W}_{\text{ar}}[0{:}V_{\text{ctc}}] \, \tilde{\mathbf{h}}_t + \mathbf{b}_{\text{ar}}[0{:}V_{\text{ctc}}].
\label{eq:tied-ctc-logits}
\end{equation}
With this design, there is no independent $\mathbf{W}_{\text{ctc}}$ parameter matrix.
Gradients from the CTC loss update the same token classifier used by the AR decoder (for the shared token subset).

\paragraph{Effect on parameter count and reporting.}
Compared to the baseline hybrid model with a separate linear CTC head, this ablation removes the dedicated CTC classifier parameters $\mathbf{W}_{\text{ctc}} \in \mathbb{R}^{V_{\text{ctc}}\times d_{\text{enc}}}$ (and optional bias), i.e., it eliminates approximately $V_{\text{ctc}}\cdot d_{\text{enc}}$ learnable weights.
No new classifier parameters are introduced because CTC reuses the AR output projection rows.
When $d_{\text{enc}}\ne d_{\text{dec}}$, we reuse the existing encoder-to-decoder projection already required for cross-attention (thus the ablation does not add an extra ``adapter'' beyond the baseline AR model).
In the thesis, we report parameter counts and MACs for each variant in a single table (baseline hybrid vs tied-CTC), and explicitly attribute any reduction in parameters to the removal of the standalone CTC head.

\paragraph{Interpretation.}
This ablation isolates whether improvements from hybrid training are attributable to the auxiliary alignment objective (CTC) rather than increased head capacity.
It also aligns with the intuition that the decoder is the primary recognizer, while CTC provides complementary supervision using the same classifier.

\subsection{Configuration switches for the ablation}
We expose minimal switches in the configuration under \texttt{dual\_head.tie}:
\begin{itemize}
  \item \texttt{ctc\_to\_ar\_outproj}: enables the tied CTC classifier (Variant A).
  \item \texttt{ar\_emb\_outproj}: optional secondary ablation that ties the AR embedding matrix to the AR output projection (LM-style weight tying).
  \item \texttt{ctc\_input\_space}: \texttt{"dec"} (recommended) uses the existing encoder$\to$decoder projection so that Eq.~\eqref{eq:tied-ctc-logits} is well-defined even when $d_{\text{enc}}\ne d_{\text{dec}}$.
\end{itemize}
A concrete YAML snippet is:
\begin{verbatim}
dual_head:
  enabled: true
  arch_ctc: linear
  lambda_ar: 1.0
  lambda_ctc: 0.6
  tie:
    ctc_to_ar_outproj: true
    ar_emb_outproj: false
    ctc_input_space: "dec"
\end{verbatim}

\subsection{Optional auxiliary supervision inside the decoder}
\label{sec:decoder-aux-supervision}
To more directly regularize the AR decoder (``between decoder layers''), we additionally implement two optional training-time losses.
These ablations are available both in hybrid AR+CTC mode (under \texttt{dual\_head.*}) and in AR-only mode (under top-level config keys), enabling fair comparisons.

\paragraph{Deep supervision (auxiliary token cross-entropy at intermediate layers).}
When enabled, we compute token-level cross-entropy not only from the final decoder layer but also from intermediate layers (e.g., layers 1--3 for a 4-layer decoder).
Crucially, all layers reuse the \emph{same} vocabulary projection matrix (the AR output projection), so this adds no independent classifier capacity.

For hybrid AR+CTC mode, the configuration lives under \texttt{dual\_head.deep\_supervision}:
\begin{verbatim}
dual_head:
  deep_supervision:
    enabled: true
    layers: [1, 2, 3]      # supervised intermediate layers (1-indexed)
    include_final: true    # also supervise the final layer
    reduce: mean           # mean|sum across supervised layers
\end{verbatim}

For AR-only mode, use a top-level block:
\begin{verbatim}
deep_supervision:
  enabled: true
  layers: [1, 2, 3]
  include_final: true
  reduce: mean
\end{verbatim}

\paragraph{Decoder-side ``CTC-like'' loss (CTC on encoder frames conditioned on decoder states).}
True CTC is defined over a frame-synchronous time axis, whereas decoder layers are naturally defined over the token axis.
To obtain a decoder-internal auxiliary signal while remaining mathematically consistent with CTC, we construct frame-length ($T'$) logits conditioned on intermediate decoder states.
The token classifier itself is still shared (we reuse the AR output projection rows for the shared vocabulary prefix), so we do not introduce an additional classifier matrix.

This option implements the suggestion to ``put the CTC head with the transformer decoder as a regularizer'' (see discussion with supervisor).

Two modes are available via \texttt{decoder\_side\_ctc.mode}:
\begin{itemize}
  \item \textbf{attention} (default): Encoder frames (queries) attend to decoder token states (keys/values) via a small reverse cross-attention module. This adds $\sim 3 d^2$ parameters (Q/K/V projections + output) but allows each frame to selectively attend to relevant decoder positions.
  \item \textbf{pool}: Global average pooling of decoder states, then broadcast-added to encoder frames. This adds \emph{zero} extra parameters and provides a simpler baseline for testing decoder-conditioning.
\end{itemize}

For hybrid AR+CTC mode, the configuration is under \texttt{dual\_head.decoder\_side\_ctc}:
\begin{verbatim}
dual_head:
  decoder_side_ctc:
    enabled: true
    mode: attention      # "attention" or "pool"
    layers: [1, 2, 3]    # which decoder layer states condition the frame logits
    lambda: 0.1          # weight of this auxiliary loss
    nhead: 4             # attention heads (only used if mode=attention)
    layernorm: true      # layer norm after attention (only if mode=attention)
\end{verbatim}

For AR-only mode (AR + decoder-side CTC ablation), use a top-level block:
\begin{verbatim}
decoder_side_ctc:
  enabled: true
  mode: pool             # zero extra params
  layers: [1, 2, 3]
  lambda: 0.1
\end{verbatim}

\paragraph{Ablation matrix.}
With these switches, the following conditions can be compared on a common encoder/decoder backbone:
\begin{enumerate}
  \item \textbf{AR-only}: pure autoregressive CE (baseline).
  \item \textbf{AR + deep supervision}: CE at layers 1/2/3/4 (no extra head capacity).
  \item \textbf{AR + decoder-side CTC}: CTC regularization ``between decoder layers'' (adds reverse-attention params only).
  \item \textbf{Hybrid AR+CTC (encoder-side)}: separate CTC head after encoder (baseline hybrid).
  \item \textbf{Hybrid AR+CTC + deep supervision / decoder-side CTC}: combinations of the above.
\end{enumerate}
This design directly tests whether placing auxiliary supervision closer to the decoder (the professor's suggestion) improves recognition compared to placing it after the encoder.

\paragraph{Secondary ablation and experimental design.}
The switch \texttt{ar\_emb\_outproj=true} enables AR embedding--classifier tying (i.e., the same weight matrix is used for token lookup and for the final vocabulary projection).
This can change optimization dynamics even without CTC weight sharing.
For clean attribution, we therefore treat it as an \emph{independent} ablation and report it both (i) on its own and (ii) in combination with Variant~A.
Concretely, we run a small factorial comparison: baseline hybrid (separate CTC head), Variant~A (CTC$\leftrightarrow$out\_proj tied), AR-only tying (\texttt{ar\_emb\_outproj=true} only), and ``full tying'' (both flags enabled).

\subsection{Results and conclusion (capacity-control test)}
\label{sec:ctc-weight-tying-results}
Table~\ref{tab:hybrid-ctc-tying-results} reports the word-level results of the AR-only baseline and hybrid AR+CTC training under different tying configurations.
All numbers are character error rate (CER) and word error rate (WER) in \%, using the same encoder/decoder and evaluation protocol across runs.
We sweep the CTC weight $\lambda_{\text{ctc}}\in\{0.5,0.6\}$, keeping $\lambda_{\text{ar}}=1$.

\begin{table}[t]
\centering
\begin{tabular}{lcc}
\hline
\textbf{Variant} & \textbf{CER (\%)} & \textbf{WER (\%)} \\
\hline
AR-only (no CTC) & 6.82 & 10.31 \\
\hline
Hybrid AR+CTC (separate CTC head), $\lambda_{\text{ctc}}=0.6$ & 6.46 & 9.76 \\
Hybrid AR+CTC (separate CTC head), $\lambda_{\text{ctc}}=0.5$ & 6.48 & 9.87 \\
\hline
Hybrid + tie \texttt{ar\_emb\_outproj} only, $\lambda_{\text{ctc}}=0.6$ & 6.53 & 9.97 \\
Hybrid + tie \texttt{ar\_emb\_outproj} only, $\lambda_{\text{ctc}}=0.5$ & 6.69 & 10.12 \\
\hline
Hybrid + tie \texttt{ctc\_to\_ar\_outproj} only, $\lambda_{\text{ctc}}=0.6$ & 7.00 & 10.49 \\
Hybrid + tie \texttt{ctc\_to\_ar\_outproj} only, $\lambda_{\text{ctc}}=0.5$ & 6.99 & 10.53 \\
\hline
Hybrid + full tying (both flags), $\lambda_{\text{ctc}}=0.5$ & 7.22 & 11.32 \\
\hline
\end{tabular}
\caption{Ablation of hybrid AR+CTC training and weight tying (word-level recognition). ``Separate CTC head'' denotes the standard hybrid model with an independent CTC classifier matrix $\mathbf{W}_{\text{ctc}}$. ``Tie \texttt{ctc\_to\_ar\_outproj}'' removes this independent classifier by reusing the AR output projection rows as in Eq.~\eqref{eq:tied-ctc-logits}.}
\label{tab:hybrid-ctc-tying-results}
\end{table}

\paragraph{What the capacity-control ablation tells us.}
The untied hybrid model (separate CTC head) improves over AR-only for both $\lambda_{\text{ctc}}=0.5$ and $0.6$.
However, when we apply the intended capacity-control by tying the CTC classifier to the AR output projection (\texttt{ctc\_to\_ar\_outproj=true}), performance drops back to (and slightly below) the AR-only baseline.
Therefore, based on these experiments, we \emph{cannot} conclude that the hybrid gains come purely from the auxiliary objective independent of added head capacity.
The results are consistent with two (not mutually exclusive) explanations: (i) the standalone CTC head provides useful classifier flexibility beyond the shared AR projection, and/or (ii) the tying constraint introduces optimization trade-offs (gradient interference) that negate the hybrid benefit.

\paragraph{Secondary observation (embedding/output tying).}
Tying the AR embedding to the AR output projection (\texttt{ar\_emb\_outproj=true}) does not improve performance in this setting and slightly degrades CER/WER compared to the untied hybrid.
This reinforces the importance of treating embedding/output tying as an independent ablation rather than as a ``free'' control.

\subsection{Reproducibility and sanity checks}
When tying is enabled, we enforce the following invariants at model initialization:
\begin{itemize}
  \item $V_{\text{ar}} \ge V_{\text{ctc}}$ so that a prefix slice exists.
  \item Character-mode special token layout: if $V_{\text{ar}} = V_{\text{ctc}} + 3$, then \texttt{PAD}$=V_{\text{ctc}}$, \texttt{BOS}$=V_{\text{ctc}}+1$, \texttt{EOS}$=V_{\text{ctc}}+2$.
  \item \texttt{ctc\_input\_space="enc"} is only allowed when $d_{\text{enc}}=d_{\text{dec}}$.
\end{itemize}
These checks ensure the tied classifier uses consistent token semantics across CTC and AR.

% End of file


\section{Hybrid autoregressive--CTC training}
\label{sec:hybrid-ar-ctc}

\paragraph{Motivation.}
Autoregressive (AR) sequence-to-sequence decoders with cross-attention are strong recognizers, but training can benefit from an auxiliary objective that encourages monotonic, frame-synchronous alignment.
Connectionist Temporal Classification (CTC) provides such a signal by defining a differentiable alignment marginalization over encoder time steps.
In this work, we adopt a \emph{hybrid} training regime where an AR decoder is the primary recognizer, while an auxiliary CTC head regularizes the shared encoder representation.

\subsection{Model}
Let $\mathbf{x} \in \mathbb{R}^{T \times C}$ denote an input IMU sequence of length $T$ with $C$ sensor channels.
An encoder $f_{\text{enc}}$ produces a downsampled sequence of hidden states
\begin{equation}
\mathbf{h} = f_{\text{enc}}(\mathbf{x}) \in \mathbb{R}^{T' \times d_{\text{enc}}},
\end{equation}
where $T' \le T$ due to temporal downsampling.
An AR decoder $f_{\text{ar}}$ attends to $\mathbf{h}$ and predicts a token sequence $\mathbf{y} = (y_1,\dots,y_N)$ via teacher forcing during training.

\paragraph{Vocabularies and special tokens.}
We use a character-level vocabulary (``categories'') with the CTC blank symbol at index $0$.
For AR decoding we use the same base vocabulary and append three special symbols \texttt{PAD}, \texttt{BOS}, \texttt{EOS} at the end.
Thus, for character mode,
\begin{equation}
V_{\text{ctc}} = |\mathcal{V}_{\text{ctc}}|, \qquad
V_{\text{ar}} = V_{\text{ctc}} + 3.
\end{equation}
This layout ensures that token IDs $0,\dots,V_{\text{ctc}}-1$ have the same semantics for both CTC and AR.

\subsection{Objectives}
\paragraph{AR cross-entropy.}
With teacher forcing, the decoder defines next-token probabilities
\begin{equation}
P_{\text{ar}}(y_t \mid y_{<t}, \mathbf{h}), \qquad t=1,\dots,N.
\end{equation}
The AR loss is the standard cross-entropy (optionally with label smoothing) summed over positions:
\begin{equation}
\mathcal{L}_{\text{ar}} = - \sum_{t=1}^{N} \log P_{\text{ar}}(y_t \mid y_{<t}, \mathbf{h}).
\end{equation}

\paragraph{CTC loss.}
A per-timestep classifier produces logits over $\mathcal{V}_{\text{ctc}}$ for each encoder frame.
CTC defines the conditional likelihood by summing over all monotonic alignments that collapse (via blank removal and repetition merging) to the target sequence.
We use the negative log-likelihood form
\begin{equation}
\mathcal{L}_{\text{ctc}} = -\log P_{\text{ctc}}(\mathbf{y} \mid \mathbf{h}).
\end{equation}

\paragraph{Hybrid loss.}
The overall hybrid loss combines both terms with weights $\lambda_{\text{ar}}$ and $\lambda_{\text{ctc}}$:
\begin{equation}
\mathcal{L}_{\text{hyb}} = \lambda_{\text{ar}}\, \mathcal{L}_{\text{ar}} + \lambda_{\text{ctc}}\, \mathcal{L}_{\text{ctc}}.
\label{eq:hybrid-loss}
\end{equation}
In our experiments, $\lambda_{\text{ar}}=1$ and $\lambda_{\text{ctc}}$ is treated as a hyperparameter (optionally scheduled across epochs).

\paragraph{Why this helps.}
The CTC term provides an auxiliary alignment/regularization signal on the encoder time axis.
This can stabilize optimization and improve generalization, especially in settings with noisy sensor sequences and variable writing styles.
Importantly, the AR decoder remains the primary recognizer at inference time.

\section{Ablation: ``single-decoder-ish'' CTC via weight tying}
\label{sec:ctc-weight-tying}

\subsection{Goal}
A common critique of hybrid models is that gains might come from adding extra parameters (the auxiliary head) rather than from the auxiliary objective.
To control for this, we introduce an ablation where the CTC classifier reuses the AR decoder's output projection weights.
This produces a more ``single-decoder-ish'' setup: CTC still scores encoder frames, but it does not introduce a separate token classifier matrix.

\subsection{Baseline (separate CTC head)}
In the standard hybrid model, CTC logits are computed via a dedicated linear layer
\begin{equation}
\mathbf{z}^{\text{ctc}}_t = \mathbf{W}_{\text{ctc}}\, \mathbf{h}_t + \mathbf{b}_{\text{ctc}},
\qquad \mathbf{W}_{\text{ctc}} \in \mathbb{R}^{V_{\text{ctc}} \times d_{\text{enc}}}.
\end{equation}
This adds $V_{\text{ctc}}\cdot d_{\text{enc}}$ parameters (plus bias).

\subsection{Variant A (primary ablation): tie CTC to the AR output projection}
\paragraph{AR output projection.}
The AR decoder produces hidden states in decoder space $\mathbf{s}_t \in \mathbb{R}^{d_{\text{dec}}}$ and maps them to vocabulary logits using its output projection
\begin{equation}
\mathbf{z}^{\text{ar}}_t = \mathbf{W}_{\text{ar}}\, \mathbf{s}_t + \mathbf{b}_{\text{ar}},
\qquad \mathbf{W}_{\text{ar}} \in \mathbb{R}^{V_{\text{ar}} \times d_{\text{dec}}}.
\end{equation}

\paragraph{CTC logits from encoder frames using shared classifier rows.}
To reuse the same classifier weights, we compute CTC logits by taking the \emph{prefix slice} of the AR projection matrix corresponding to the shared vocabulary IDs:
\begin{equation}
\mathbf{W}_{\text{ctc}} \equiv \mathbf{W}_{\text{ar}}[0{:}V_{\text{ctc}}],
\qquad \mathbf{b}_{\text{ctc}} \equiv \mathbf{b}_{\text{ar}}[0{:}V_{\text{ctc}}].
\end{equation}
Since CTC operates on encoder frames, we must ensure its input feature dimension matches $d_{\text{dec}}$.
When $d_{\text{enc}} \ne d_{\text{dec}}$, we reuse the existing encoder-to-decoder projection (already required for AR cross-attention in our implementation):
\begin{equation}
\tilde{\mathbf{h}}_t = \mathbf{P}\, \mathbf{h}_t, \qquad \mathbf{P} \in \mathbb{R}^{d_{\text{dec}} \times d_{\text{enc}}}.
\end{equation}
The tied CTC logits become
\begin{equation}
\mathbf{z}^{\text{ctc}}_t = \mathbf{W}_{\text{ar}}[0{:}V_{\text{ctc}}] \, \tilde{\mathbf{h}}_t + \mathbf{b}_{\text{ar}}[0{:}V_{\text{ctc}}].
\label{eq:tied-ctc-logits}
\end{equation}
With this design, there is no independent $\mathbf{W}_{\text{ctc}}$ parameter matrix.
Gradients from the CTC loss update the same token classifier used by the AR decoder (for the shared token subset).

\paragraph{Effect on parameter count and reporting.}
Compared to the baseline hybrid model with a separate linear CTC head, this ablation removes the dedicated CTC classifier parameters $\mathbf{W}_{\text{ctc}} \in \mathbb{R}^{V_{\text{ctc}}\times d_{\text{enc}}}$ (and optional bias), i.e., it eliminates approximately $V_{\text{ctc}}\cdot d_{\text{enc}}$ learnable weights.
No new classifier parameters are introduced because CTC reuses the AR output projection rows.
When $d_{\text{enc}}\ne d_{\text{dec}}$, we reuse the existing encoder-to-decoder projection already required for cross-attention (thus the ablation does not add an extra ``adapter'' beyond the baseline AR model).
In the thesis, we report parameter counts and MACs for each variant in a single table (baseline hybrid vs tied-CTC), and explicitly attribute any reduction in parameters to the removal of the standalone CTC head.

\paragraph{Interpretation.}
This ablation isolates whether improvements from hybrid training are attributable to the auxiliary alignment objective (CTC) rather than increased head capacity.
It also aligns with the intuition that the decoder is the primary recognizer, while CTC provides complementary supervision using the same classifier.

\subsection{Configuration switches for the ablation}
We expose minimal switches in the configuration under \texttt{dual\_head.tie}:
\begin{itemize}
  \item \texttt{ctc\_to\_ar\_outproj}: enables the tied CTC classifier (Variant A).
  \item \texttt{ar\_emb\_outproj}: optional secondary ablation that ties the AR embedding matrix to the AR output projection (LM-style weight tying).
  \item \texttt{ctc\_input\_space}: \texttt{"dec"} (recommended) uses the existing encoder$\to$decoder projection so that Eq.~\eqref{eq:tied-ctc-logits} is well-defined even when $d_{\text{enc}}\ne d_{\text{dec}}$.
\end{itemize}
A concrete YAML snippet is:
\begin{verbatim}
dual_head:
  enabled: true
  arch_ctc: linear
  lambda_ar: 1.0
  lambda_ctc: 0.6
  tie:
    ctc_to_ar_outproj: true
    ar_emb_outproj: false
    ctc_input_space: "dec"
\end{verbatim}

\subsection{Optional auxiliary supervision inside the decoder}
\label{sec:decoder-aux-supervision}
To more directly regularize the AR decoder (``between decoder layers''), we additionally implement two optional training-time losses.
These ablations are available both in hybrid AR+CTC mode (under \texttt{dual\_head.*}) and in AR-only mode (under top-level config keys), enabling fair comparisons.

\paragraph{Deep supervision (auxiliary token cross-entropy at intermediate layers).}
When enabled, we compute token-level cross-entropy not only from the final decoder layer but also from intermediate layers (e.g., layers 1--3 for a 4-layer decoder).
Crucially, all layers reuse the \emph{same} vocabulary projection matrix (the AR output projection), so this adds no independent classifier capacity.

For hybrid AR+CTC mode, the configuration lives under \texttt{dual\_head.deep\_supervision}:
\begin{verbatim}
dual_head:
  deep_supervision:
    enabled: true
    layers: [1, 2, 3]      # supervised intermediate layers (1-indexed)
    include_final: true    # also supervise the final layer
    reduce: mean           # mean|sum across supervised layers
\end{verbatim}

For AR-only mode, use a top-level block:
\begin{verbatim}
deep_supervision:
  enabled: true
  layers: [1, 2, 3]
  include_final: true
  reduce: mean
\end{verbatim}

\paragraph{Decoder-side ``CTC-like'' loss (CTC on encoder frames conditioned on decoder states).}
True CTC is defined over a frame-synchronous time axis, whereas decoder layers are naturally defined over the token axis.
To obtain a decoder-internal auxiliary signal while remaining mathematically consistent with CTC, we construct frame-length ($T'$) logits conditioned on intermediate decoder states.
The token classifier itself is still shared (we reuse the AR output projection rows for the shared vocabulary prefix), so we do not introduce an additional classifier matrix.

This option implements the suggestion to ``put the CTC head with the transformer decoder as a regularizer'' (see discussion with supervisor).

Two modes are available via \texttt{decoder\_side\_ctc.mode}:
\begin{itemize}
  \item \textbf{attention} (default): Encoder frames (queries) attend to decoder token states (keys/values) via a small reverse cross-attention module. This adds $\sim 3 d^2$ parameters (Q/K/V projections + output) but allows each frame to selectively attend to relevant decoder positions.
  \item \textbf{pool}: Global average pooling of decoder states, then broadcast-added to encoder frames. This adds \emph{zero} extra parameters and provides a simpler baseline for testing decoder-conditioning.
\end{itemize}

For hybrid AR+CTC mode, the configuration is under \texttt{dual\_head.decoder\_side\_ctc}:
\begin{verbatim}
dual_head:
  decoder_side_ctc:
    enabled: true
    mode: attention      # "attention" or "pool"
    layers: [1, 2, 3]    # which decoder layer states condition the frame logits
    lambda: 0.1          # weight of this auxiliary loss
    nhead: 4             # attention heads (only used if mode=attention)
    layernorm: true      # layer norm after attention (only if mode=attention)
\end{verbatim}

For AR-only mode (AR + decoder-side CTC ablation), use a top-level block:
\begin{verbatim}
decoder_side_ctc:
  enabled: true
  mode: pool             # zero extra params
  layers: [1, 2, 3]
  lambda: 0.1
\end{verbatim}

\paragraph{Ablation matrix.}
With these switches, the following conditions can be compared on a common encoder/decoder backbone:
\begin{enumerate}
  \item \textbf{AR-only}: pure autoregressive CE (baseline).
  \item \textbf{AR + deep supervision}: CE at layers 1/2/3/4 (no extra head capacity).
  \item \textbf{AR + decoder-side CTC}: CTC regularization ``between decoder layers'' (adds reverse-attention params only).
  \item \textbf{Hybrid AR+CTC (encoder-side)}: separate CTC head after encoder (baseline hybrid).
  \item \textbf{Hybrid AR+CTC + deep supervision / decoder-side CTC}: combinations of the above.
\end{enumerate}
This design directly tests whether placing auxiliary supervision closer to the decoder (the professor's suggestion) improves recognition compared to placing it after the encoder.

\paragraph{Secondary ablation and experimental design.}
The switch \texttt{ar\_emb\_outproj=true} enables AR embedding--classifier tying (i.e., the same weight matrix is used for token lookup and for the final vocabulary projection).
This can change optimization dynamics even without CTC weight sharing.
For clean attribution, we therefore treat it as an \emph{independent} ablation and report it both (i) on its own and (ii) in combination with Variant~A.
Concretely, we run a small factorial comparison: baseline hybrid (separate CTC head), Variant~A (CTC$\leftrightarrow$out\_proj tied), AR-only tying (\texttt{ar\_emb\_outproj=true} only), and ``full tying'' (both flags enabled).

\subsection{Results and conclusion (capacity-control test)}
\label{sec:ctc-weight-tying-results}
Table~\ref{tab:hybrid-ctc-tying-results} reports the word-level results of the AR-only baseline and hybrid AR+CTC training under different tying configurations.
All numbers are character error rate (CER) and word error rate (WER) in \%, using the same encoder/decoder and evaluation protocol across runs.
We sweep the CTC weight $\lambda_{\text{ctc}}\in\{0.5,0.6\}$, keeping $\lambda_{\text{ar}}=1$.

\begin{table}[t]
\centering
\begin{tabular}{lcc}
\hline
\textbf{Variant} & \textbf{CER (\%)} & \textbf{WER (\%)} \\
\hline
AR-only (no CTC) & 6.82 & 10.31 \\
\hline
Hybrid AR+CTC (separate CTC head), $\lambda_{\text{ctc}}=0.6$ & 6.46 & 9.76 \\
Hybrid AR+CTC (separate CTC head), $\lambda_{\text{ctc}}=0.5$ & 6.48 & 9.87 \\
\hline
Hybrid + tie \texttt{ar\_emb\_outproj} only, $\lambda_{\text{ctc}}=0.6$ & 6.53 & 9.97 \\
Hybrid + tie \texttt{ar\_emb\_outproj} only, $\lambda_{\text{ctc}}=0.5$ & 6.69 & 10.12 \\
\hline
Hybrid + tie \texttt{ctc\_to\_ar\_outproj} only, $\lambda_{\text{ctc}}=0.6$ & 7.00 & 10.49 \\
Hybrid + tie \texttt{ctc\_to\_ar\_outproj} only, $\lambda_{\text{ctc}}=0.5$ & 6.99 & 10.53 \\
\hline
Hybrid + full tying (both flags), $\lambda_{\text{ctc}}=0.5$ & 7.22 & 11.32 \\
\hline
\end{tabular}
\caption{Ablation of hybrid AR+CTC training and weight tying (word-level recognition). ``Separate CTC head'' denotes the standard hybrid model with an independent CTC classifier matrix $\mathbf{W}_{\text{ctc}}$. ``Tie \texttt{ctc\_to\_ar\_outproj}'' removes this independent classifier by reusing the AR output projection rows as in Eq.~\eqref{eq:tied-ctc-logits}.}
\label{tab:hybrid-ctc-tying-results}
\end{table}

\paragraph{What the capacity-control ablation tells us.}
The untied hybrid model (separate CTC head) improves over AR-only for both $\lambda_{\text{ctc}}=0.5$ and $0.6$.
However, when we apply the intended capacity-control by tying the CTC classifier to the AR output projection (\texttt{ctc\_to\_ar\_outproj=true}), performance drops back to (and slightly below) the AR-only baseline.
Therefore, based on these experiments, we \emph{cannot} conclude that the hybrid gains come purely from the auxiliary objective independent of added head capacity.
The results are consistent with two (not mutually exclusive) explanations: (i) the standalone CTC head provides useful classifier flexibility beyond the shared AR projection, and/or (ii) the tying constraint introduces optimization trade-offs (gradient interference) that negate the hybrid benefit.

\paragraph{Secondary observation (embedding/output tying).}
Tying the AR embedding to the AR output projection (\texttt{ar\_emb\_outproj=true}) does not improve performance in this setting and slightly degrades CER/WER compared to the untied hybrid.
This reinforces the importance of treating embedding/output tying as an independent ablation rather than as a ``free'' control.

\subsection{Reproducibility and sanity checks}
When tying is enabled, we enforce the following invariants at model initialization:
\begin{itemize}
  \item $V_{\text{ar}} \ge V_{\text{ctc}}$ so that a prefix slice exists.
  \item Character-mode special token layout: if $V_{\text{ar}} = V_{\text{ctc}} + 3$, then \texttt{PAD}$=V_{\text{ctc}}$, \texttt{BOS}$=V_{\text{ctc}}+1$, \texttt{EOS}$=V_{\text{ctc}}+2$.
  \item \texttt{ctc\_input\_space="enc"} is only allowed when $d_{\text{enc}}=d_{\text{dec}}$.
\end{itemize}
These checks ensure the tied classifier uses consistent token semantics across CTC and AR.

% End of file


\section{Hybrid autoregressive--CTC training}
\label{sec:hybrid-ar-ctc}

\paragraph{Motivation.}
Autoregressive (AR) sequence-to-sequence decoders with cross-attention are strong recognizers, but training can benefit from an auxiliary objective that encourages monotonic, frame-synchronous alignment.
Connectionist Temporal Classification (CTC) provides such a signal by defining a differentiable alignment marginalization over encoder time steps.
In this work, we adopt a \emph{hybrid} training regime where an AR decoder is the primary recognizer, while an auxiliary CTC head regularizes the shared encoder representation.

\subsection{Model}
Let $\mathbf{x} \in \mathbb{R}^{T \times C}$ denote an input IMU sequence of length $T$ with $C$ sensor channels.
An encoder $f_{\text{enc}}$ produces a downsampled sequence of hidden states
\begin{equation}
\mathbf{h} = f_{\text{enc}}(\mathbf{x}) \in \mathbb{R}^{T' \times d_{\text{enc}}},
\end{equation}
where $T' \le T$ due to temporal downsampling.
An AR decoder $f_{\text{ar}}$ attends to $\mathbf{h}$ and predicts a token sequence $\mathbf{y} = (y_1,\dots,y_N)$ via teacher forcing during training.

\paragraph{Vocabularies and special tokens.}
We use a character-level vocabulary (``categories'') with the CTC blank symbol at index $0$.
For AR decoding we use the same base vocabulary and append three special symbols \texttt{PAD}, \texttt{BOS}, \texttt{EOS} at the end.
Thus, for character mode,
\begin{equation}
V_{\text{ctc}} = |\mathcal{V}_{\text{ctc}}|, \qquad
V_{\text{ar}} = V_{\text{ctc}} + 3.
\end{equation}
This layout ensures that token IDs $0,\dots,V_{\text{ctc}}-1$ have the same semantics for both CTC and AR.

\subsection{Objectives}
\paragraph{AR cross-entropy.}
With teacher forcing, the decoder defines next-token probabilities
\begin{equation}
P_{\text{ar}}(y_t \mid y_{<t}, \mathbf{h}), \qquad t=1,\dots,N.
\end{equation}
The AR loss is the standard cross-entropy (optionally with label smoothing) summed over positions:
\begin{equation}
\mathcal{L}_{\text{ar}} = - \sum_{t=1}^{N} \log P_{\text{ar}}(y_t \mid y_{<t}, \mathbf{h}).
\end{equation}

\paragraph{CTC loss.}
A per-timestep classifier produces logits over $\mathcal{V}_{\text{ctc}}$ for each encoder frame.
CTC defines the conditional likelihood by summing over all monotonic alignments that collapse (via blank removal and repetition merging) to the target sequence.
We use the negative log-likelihood form
\begin{equation}
\mathcal{L}_{\text{ctc}} = -\log P_{\text{ctc}}(\mathbf{y} \mid \mathbf{h}).
\end{equation}

\paragraph{Hybrid loss.}
The overall hybrid loss combines both terms with weights $\lambda_{\text{ar}}$ and $\lambda_{\text{ctc}}$:
\begin{equation}
\mathcal{L}_{\text{hyb}} = \lambda_{\text{ar}}\, \mathcal{L}_{\text{ar}} + \lambda_{\text{ctc}}\, \mathcal{L}_{\text{ctc}}.
\label{eq:hybrid-loss}
\end{equation}
In our experiments, $\lambda_{\text{ar}}=1$ and $\lambda_{\text{ctc}}$ is treated as a hyperparameter (optionally scheduled across epochs).

\paragraph{Why this helps.}
The CTC term provides an auxiliary alignment/regularization signal on the encoder time axis.
This can stabilize optimization and improve generalization, especially in settings with noisy sensor sequences and variable writing styles.
Importantly, the AR decoder remains the primary recognizer at inference time.

\section{Ablation: ``single-decoder-ish'' CTC via weight tying}
\label{sec:ctc-weight-tying}

\subsection{Goal}
A common critique of hybrid models is that gains might come from adding extra parameters (the auxiliary head) rather than from the auxiliary objective.
To control for this, we introduce an ablation where the CTC classifier reuses the AR decoder's output projection weights.
This produces a more ``single-decoder-ish'' setup: CTC still scores encoder frames, but it does not introduce a separate token classifier matrix.

\subsection{Baseline (separate CTC head)}
In the standard hybrid model, CTC logits are computed via a dedicated linear layer
\begin{equation}
\mathbf{z}^{\text{ctc}}_t = \mathbf{W}_{\text{ctc}}\, \mathbf{h}_t + \mathbf{b}_{\text{ctc}},
\qquad \mathbf{W}_{\text{ctc}} \in \mathbb{R}^{V_{\text{ctc}} \times d_{\text{enc}}}.
\end{equation}
This adds $V_{\text{ctc}}\cdot d_{\text{enc}}$ parameters (plus bias).

\subsection{Variant A (primary ablation): tie CTC to the AR output projection}
\paragraph{AR output projection.}
The AR decoder produces hidden states in decoder space $\mathbf{s}_t \in \mathbb{R}^{d_{\text{dec}}}$ and maps them to vocabulary logits using its output projection
\begin{equation}
\mathbf{z}^{\text{ar}}_t = \mathbf{W}_{\text{ar}}\, \mathbf{s}_t + \mathbf{b}_{\text{ar}},
\qquad \mathbf{W}_{\text{ar}} \in \mathbb{R}^{V_{\text{ar}} \times d_{\text{dec}}}.
\end{equation}

\paragraph{CTC logits from encoder frames using shared classifier rows.}
To reuse the same classifier weights, we compute CTC logits by taking the \emph{prefix slice} of the AR projection matrix corresponding to the shared vocabulary IDs:
\begin{equation}
\mathbf{W}_{\text{ctc}} \equiv \mathbf{W}_{\text{ar}}[0{:}V_{\text{ctc}}],
\qquad \mathbf{b}_{\text{ctc}} \equiv \mathbf{b}_{\text{ar}}[0{:}V_{\text{ctc}}].
\end{equation}
Since CTC operates on encoder frames, we must ensure its input feature dimension matches $d_{\text{dec}}$.
When $d_{\text{enc}} \ne d_{\text{dec}}$, we reuse the existing encoder-to-decoder projection (already required for AR cross-attention in our implementation):
\begin{equation}
\tilde{\mathbf{h}}_t = \mathbf{P}\, \mathbf{h}_t, \qquad \mathbf{P} \in \mathbb{R}^{d_{\text{dec}} \times d_{\text{enc}}}.
\end{equation}
The tied CTC logits become
\begin{equation}
\mathbf{z}^{\text{ctc}}_t = \mathbf{W}_{\text{ar}}[0{:}V_{\text{ctc}}] \, \tilde{\mathbf{h}}_t + \mathbf{b}_{\text{ar}}[0{:}V_{\text{ctc}}].
\label{eq:tied-ctc-logits}
\end{equation}
With this design, there is no independent $\mathbf{W}_{\text{ctc}}$ parameter matrix.
Gradients from the CTC loss update the same token classifier used by the AR decoder (for the shared token subset).

\paragraph{Effect on parameter count and reporting.}
Compared to the baseline hybrid model with a separate linear CTC head, this ablation removes the dedicated CTC classifier parameters $\mathbf{W}_{\text{ctc}} \in \mathbb{R}^{V_{\text{ctc}}\times d_{\text{enc}}}$ (and optional bias), i.e., it eliminates approximately $V_{\text{ctc}}\cdot d_{\text{enc}}$ learnable weights.
No new classifier parameters are introduced because CTC reuses the AR output projection rows.
When $d_{\text{enc}}\ne d_{\text{dec}}$, we reuse the existing encoder-to-decoder projection already required for cross-attention (thus the ablation does not add an extra ``adapter'' beyond the baseline AR model).
In the thesis, we report parameter counts and MACs for each variant in a single table (baseline hybrid vs tied-CTC), and explicitly attribute any reduction in parameters to the removal of the standalone CTC head.

\paragraph{Interpretation.}
This ablation isolates whether improvements from hybrid training are attributable to the auxiliary alignment objective (CTC) rather than increased head capacity.
It also aligns with the intuition that the decoder is the primary recognizer, while CTC provides complementary supervision using the same classifier.

\subsection{Configuration switches for the ablation}
We expose minimal switches in the configuration under \texttt{dual\_head.tie}:
\begin{itemize}
  \item \texttt{ctc\_to\_ar\_outproj}: enables the tied CTC classifier (Variant A).
  \item \texttt{ar\_emb\_outproj}: optional secondary ablation that ties the AR embedding matrix to the AR output projection (LM-style weight tying).
  \item \texttt{ctc\_input\_space}: \texttt{"dec"} (recommended) uses the existing encoder$\to$decoder projection so that Eq.~\eqref{eq:tied-ctc-logits} is well-defined even when $d_{\text{enc}}\ne d_{\text{dec}}$.
\end{itemize}
A concrete YAML snippet is:
\begin{verbatim}
dual_head:
  enabled: true
  arch_ctc: linear
  lambda_ar: 1.0
  lambda_ctc: 0.6
  tie:
    ctc_to_ar_outproj: true
    ar_emb_outproj: false
    ctc_input_space: "dec"
\end{verbatim}

\subsection{Optional auxiliary supervision inside the decoder}
\label{sec:decoder-aux-supervision}
To more directly regularize the AR decoder (``between decoder layers''), we additionally implement two optional training-time losses.
These ablations are available both in hybrid AR+CTC mode (under \texttt{dual\_head.*}) and in AR-only mode (under top-level config keys), enabling fair comparisons.

\paragraph{Deep supervision (auxiliary token cross-entropy at intermediate layers).}
When enabled, we compute token-level cross-entropy not only from the final decoder layer but also from intermediate layers (e.g., layers 1--3 for a 4-layer decoder).
Crucially, all layers reuse the \emph{same} vocabulary projection matrix (the AR output projection), so this adds no independent classifier capacity.

For hybrid AR+CTC mode, the configuration lives under \texttt{dual\_head.deep\_supervision}:
\begin{verbatim}
dual_head:
  deep_supervision:
    enabled: true
    layers: [1, 2, 3]      # supervised intermediate layers (1-indexed)
    include_final: true    # also supervise the final layer
    reduce: mean           # mean|sum across supervised layers
\end{verbatim}

For AR-only mode, use a top-level block:
\begin{verbatim}
deep_supervision:
  enabled: true
  layers: [1, 2, 3]
  include_final: true
  reduce: mean
\end{verbatim}

\paragraph{Decoder-side ``CTC-like'' loss (CTC on encoder frames conditioned on decoder states).}
True CTC is defined over a frame-synchronous time axis, whereas decoder layers are naturally defined over the token axis.
To obtain a decoder-internal auxiliary signal while remaining mathematically consistent with CTC, we construct frame-length ($T'$) logits conditioned on intermediate decoder states.
The token classifier itself is still shared (we reuse the AR output projection rows for the shared vocabulary prefix), so we do not introduce an additional classifier matrix.

This option implements the suggestion to ``put the CTC head with the transformer decoder as a regularizer'' (see discussion with supervisor).

Two modes are available via \texttt{decoder\_side\_ctc.mode}:
\begin{itemize}
  \item \textbf{attention} (default): Encoder frames (queries) attend to decoder token states (keys/values) via a small reverse cross-attention module. This adds $\sim 3 d^2$ parameters (Q/K/V projections + output) but allows each frame to selectively attend to relevant decoder positions.
  \item \textbf{pool}: Global average pooling of decoder states, then broadcast-added to encoder frames. This adds \emph{zero} extra parameters and provides a simpler baseline for testing decoder-conditioning.
\end{itemize}

For hybrid AR+CTC mode, the configuration is under \texttt{dual\_head.decoder\_side\_ctc}:
\begin{verbatim}
dual_head:
  decoder_side_ctc:
    enabled: true
    mode: attention      # "attention" or "pool"
    layers: [1, 2, 3]    # which decoder layer states condition the frame logits
    lambda: 0.1          # weight of this auxiliary loss
    nhead: 4             # attention heads (only used if mode=attention)
    layernorm: true      # layer norm after attention (only if mode=attention)
\end{verbatim}

For AR-only mode (AR + decoder-side CTC ablation), use a top-level block:
\begin{verbatim}
decoder_side_ctc:
  enabled: true
  mode: pool             # zero extra params
  layers: [1, 2, 3]
  lambda: 0.1
\end{verbatim}

\paragraph{Ablation matrix.}
With these switches, the following conditions can be compared on a common encoder/decoder backbone:
\begin{enumerate}
  \item \textbf{AR-only}: pure autoregressive CE (baseline).
  \item \textbf{AR + deep supervision}: CE at layers 1/2/3/4 (no extra head capacity).
  \item \textbf{AR + decoder-side CTC}: CTC regularization ``between decoder layers'' (adds reverse-attention params only).
  \item \textbf{Hybrid AR+CTC (encoder-side)}: separate CTC head after encoder (baseline hybrid).
  \item \textbf{Hybrid AR+CTC + deep supervision / decoder-side CTC}: combinations of the above.
\end{enumerate}
This design directly tests whether placing auxiliary supervision closer to the decoder (the professor's suggestion) improves recognition compared to placing it after the encoder.

\paragraph{Secondary ablation and experimental design.}
The switch \texttt{ar\_emb\_outproj=true} enables AR embedding--classifier tying (i.e., the same weight matrix is used for token lookup and for the final vocabulary projection).
This can change optimization dynamics even without CTC weight sharing.
For clean attribution, we therefore treat it as an \emph{independent} ablation and report it both (i) on its own and (ii) in combination with Variant~A.
Concretely, we run a small factorial comparison: baseline hybrid (separate CTC head), Variant~A (CTC$\leftrightarrow$out\_proj tied), AR-only tying (\texttt{ar\_emb\_outproj=true} only), and ``full tying'' (both flags enabled).

\subsection{Results and conclusion (capacity-control test)}
\label{sec:ctc-weight-tying-results}
Table~\ref{tab:hybrid-ctc-tying-results} reports the word-level results of the AR-only baseline and hybrid AR+CTC training under different tying configurations.
All numbers are character error rate (CER) and word error rate (WER) in \%, using the same encoder/decoder and evaluation protocol across runs.
We sweep the CTC weight $\lambda_{\text{ctc}}\in\{0.5,0.6\}$, keeping $\lambda_{\text{ar}}=1$.

\begin{table}[t]
\centering
\begin{tabular}{lcc}
\hline
\textbf{Variant} & \textbf{CER (\%)} & \textbf{WER (\%)} \\
\hline
AR-only (no CTC) & 6.82 & 10.31 \\
\hline
Hybrid AR+CTC (separate CTC head), $\lambda_{\text{ctc}}=0.6$ & 6.46 & 9.76 \\
Hybrid AR+CTC (separate CTC head), $\lambda_{\text{ctc}}=0.5$ & 6.48 & 9.87 \\
\hline
Hybrid + tie \texttt{ar\_emb\_outproj} only, $\lambda_{\text{ctc}}=0.6$ & 6.53 & 9.97 \\
Hybrid + tie \texttt{ar\_emb\_outproj} only, $\lambda_{\text{ctc}}=0.5$ & 6.69 & 10.12 \\
\hline
Hybrid + tie \texttt{ctc\_to\_ar\_outproj} only, $\lambda_{\text{ctc}}=0.6$ & 7.00 & 10.49 \\
Hybrid + tie \texttt{ctc\_to\_ar\_outproj} only, $\lambda_{\text{ctc}}=0.5$ & 6.99 & 10.53 \\
\hline
Hybrid + full tying (both flags), $\lambda_{\text{ctc}}=0.5$ & 7.22 & 11.32 \\
\hline
\end{tabular}
\caption{Ablation of hybrid AR+CTC training and weight tying (word-level recognition). ``Separate CTC head'' denotes the standard hybrid model with an independent CTC classifier matrix $\mathbf{W}_{\text{ctc}}$. ``Tie \texttt{ctc\_to\_ar\_outproj}'' removes this independent classifier by reusing the AR output projection rows as in Eq.~\eqref{eq:tied-ctc-logits}.}
\label{tab:hybrid-ctc-tying-results}
\end{table}

\paragraph{What the capacity-control ablation tells us.}
The untied hybrid model (separate CTC head) improves over AR-only for both $\lambda_{\text{ctc}}=0.5$ and $0.6$.
However, when we apply the intended capacity-control by tying the CTC classifier to the AR output projection (\texttt{ctc\_to\_ar\_outproj=true}), performance drops back to (and slightly below) the AR-only baseline.
Therefore, based on these experiments, we \emph{cannot} conclude that the hybrid gains come purely from the auxiliary objective independent of added head capacity.
The results are consistent with two (not mutually exclusive) explanations: (i) the standalone CTC head provides useful classifier flexibility beyond the shared AR projection, and/or (ii) the tying constraint introduces optimization trade-offs (gradient interference) that negate the hybrid benefit.

\paragraph{Secondary observation (embedding/output tying).}
Tying the AR embedding to the AR output projection (\texttt{ar\_emb\_outproj=true}) does not improve performance in this setting and slightly degrades CER/WER compared to the untied hybrid.
This reinforces the importance of treating embedding/output tying as an independent ablation rather than as a ``free'' control.

\subsection{Reproducibility and sanity checks}
When tying is enabled, we enforce the following invariants at model initialization:
\begin{itemize}
  \item $V_{\text{ar}} \ge V_{\text{ctc}}$ so that a prefix slice exists.
  \item Character-mode special token layout: if $V_{\text{ar}} = V_{\text{ctc}} + 3$, then \texttt{PAD}$=V_{\text{ctc}}$, \texttt{BOS}$=V_{\text{ctc}}+1$, \texttt{EOS}$=V_{\text{ctc}}+2$.
  \item \texttt{ctc\_input\_space="enc"} is only allowed when $d_{\text{enc}}=d_{\text{dec}}$.
\end{itemize}
These checks ensure the tied classifier uses consistent token semantics across CTC and AR.

% End of file
